\documentclass[sigconf,nonacm]{acmart}
\settopmatter{printacmref=false}
\setcopyright{none}
\renewcommand\footnotetextcopyrightpermission[1]{}
\pagestyle{plain}

\usepackage{booktabs}

\begin{document}

\title{DistributedPanGA Objective O1: HySortK-Inspired Extraction}
\subtitle{Final report}

\author{Harshaan Chugh}
\affiliation{%
  \institution{Cornell University}
  \city{Ithaca}
  \state{NY}
  \country{USA}}

\begin{abstract}
DistributedPanGA distributes FastGA's all-vs-all pangenome alignment over MPI.
Objective O1 of our proposal was to speed up its first stage, seed extraction,
using ideas from the HySortK distributed k-mer counter.
Measurement showed that extraction was dominated not by sorting or
communication but by one \texttt{lseek} system call per extracted record.
Packing each seed from the sequence window already in memory (O1a) made
extraction 34--136$\times$ faster and the whole pipeline 1.7--3.7$\times$
faster on one node. Every run produced output identical to FastGA's.
After O1a, extraction is about 2\% of the run time. The remaining planned
HySortK adaptations (threaded extraction, destination-aware staging and
overlapped exchange) could not pay off and were dropped. Supermer compression
does not apply to our seeds. We also report which HySortK ideas transfer to
this pipeline and which do not.
\end{abstract}

\maketitle

\section{Introduction}

FastGA~\cite{fastga} aligns genomes on a single node. DistributedPanGA runs
FastGA's self-alignment across MPI ranks in six stages:
(1)~\emph{extract} 40-mer seeds chosen by closed (12,8)-syncmers, keeping
each occurrence's contig, position and strand;
(2)~\emph{sort} them globally by a distributed sort with a prefix-range
exchange;
(3)~\emph{merge} equal k-mers into seed pairs with FastGA's adaptive-seed
merge;
(4)~\emph{redistribute} seeds to the rank that owns each contig pair;
(5)~\emph{align} each contig pair; and
(6)~\emph{convert} the result to PAF.

HySortK~\cite{hysortk} is a distributed sorting-based k-mer counter. It uses
supermers, a task layer with balanced dispatch, hybrid MPI/OpenMP and batched
non-blocking exchange. Objective~O1 was to adapt its ideas to stages~1--2,
while keeping every occurrence and alignments identical to FastGA's.

\section{Verification method}

Every claim below is checked against an \emph{independent} reference: FastGA
(revision \texttt{178afcb}) run on the same input. The pipeline's per-rank PAF
files are concatenated and sorted without removing duplicates, then compared
byte for byte with FastGA's PAF.

The main dataset is yeast7, which reproduces the historical 43,126
alignments. Additional phase verifiers check extraction (byte-identical to
FastGA's index), the local and distributed sorts, and the LCP values.

Timings use one Perlmutter CPU node (2$\times$AMD EPYC 7763) with 16 cores per
rank, one thread per rank, and 3 interleaved repetitions. Reported spreads are
sample standard deviations.

\section{Diagnosis (O1-0)}

Per-rank instrumentation on yeast7 (job 59679320) found:

\begin{itemize}
\item \textbf{Work was balanced.} Records per rank were within $\pm$11\% at 8
ranks.
\item \textbf{The cost per record grew with the rank count}: 0.78, 0.95, 1.48
and 3.97\,$\mu$s at 1, 2, 4 and 8 ranks.
\item \textbf{User time scaled as $1/N$, but kernel time per rank stayed
flat} at 19--30\,s. Total kernel time rose from 26\,s to 238\,s for the same
total work.
\item \textbf{One \texttt{lseek} per record.} \texttt{strace} on the serial
extractor counted 38,035 \texttt{lseek} calls for 38,023 records. Every seed
re-read 40 bases through \texttt{Get\_Contig\_Piece}, whose \texttt{fseeko}
issues a seek.
\item \textbf{Output does not depend on pre-sort order.} Shuffling the
received records before sorting left the per-rank PAFs byte-identical (yeast7,
1 and 4 ranks, 3 seeds). This is evidence for this input, not a proof.
\end{itemize}

\section{O1a: pack from the scan window}

The scanner already holds each contig window in memory. O1a packs each seed's
forward and reverse-complement k-mer straight from that window.

Windows after the first re-read 40 bases of left context, so seeds near a
window's start are covered. A forward k-mer that crosses a non-final window's
right edge falls back to the old read. On these datasets every contig fits in
one window, so the fallback never fired.

Seed selection, tie handling, coordinates, emission order and LCP handling
are unchanged.

\begin{table}[t]
\caption{O1a on yeast7, same node, 3 interleaved repetitions (job 59691033).
Kernel time is summed over ranks.}
\label{tab:o1a}
\small
\begin{tabular}{@{}lrrrr@{}}
\toprule
 & \multicolumn{2}{c}{1 rank} & \multicolumn{2}{c}{8 ranks} \\
 & before & O1a & before & O1a \\
\midrule
wall time (s) & 113.80 & \textbf{65.43} & 39.11 & \textbf{10.57} \\
extraction (s) & 49.76 & \textbf{1.48} & 28.61 & \textbf{0.21} \\
kernel time (s) & 26.7 & 2.7 & 199.3 & 6.3 \\
alignments vs.\ gold & match & match & match & match \\
\bottomrule
\end{tabular}
\end{table}

Table~\ref{tab:o1a} shows extraction 34$\times$ faster at 1 rank and
136$\times$ faster at 8 ranks. The 1-to-8-rank speedup of the whole run rose
from 2.9$\times$ to 6.2$\times$.

Removing exactly the per-record seeks removed 97\% of the 1-rank extraction
time, which strongly supports the diagnosis. Which kernel path was contended
at 8 ranks remains unattributed; it no longer affects performance.

All 20 O1a timing runs matched gold, as did all phase verifiers and an
ASan/UBSan build of the multi-window extraction tests.

\section{Dropped steps and why}

\begin{table*}[t]
\caption{Planned O1 steps that were not built, with the measurement that
decided each.}
\label{tab:dropped}
\small
\begin{tabular}{@{}p{0.2\textwidth}p{0.75\textwidth}@{}}
\toprule
Step & Decision and evidence \\
\midrule
O1b: threaded extraction &
Dropped. After O1a, extraction is 0.21 of 10.57\,s at 8 ranks ($\approx$2\%).
Threads could save at most $\approx$1.5\,s at 1 rank. \\
O1c: destination-aware staging (HySortK's per-thread $\times$
per-destination buffers) &
Dropped. The speed gain is bounded by the histogram and packing passes,
$\approx$0.45\,s. When proposed it could not lower peak memory either, because
the peak was in the alignment stage. Once later work freed those buffers,
stage~2 became the 1-rank peak (2.96\,GiB), so a memory-only variant is a
possible future option. \\
O1d: batched, overlapped exchange &
Dropped. The k-mer exchange takes $\le$0.13\,s per rank on one node, so
there is nothing to hide. Exchanges above 2\,GB per rank were later handled
with bounded rounds instead. \\
Supermer compression &
Not applicable. A seed's destination is its own first 5 bases, and
consecutive selected seeds ($\approx$2.5 bases apart) reach the same rank
with probability $\approx 1/P$. The expected run is $P/(P{-}1)$ seeds, a net
loss after headers. Routing by minimizer would split the 12-base prefix
groups that the merge needs kept together. \\
\bottomrule
\end{tabular}
\end{table*}

Table~\ref{tab:dropped} lists the other planned sub-steps. Each was evaluated
against measurements taken after O1a.

\section{What transfers from HySortK}

\begin{itemize}
\item \textbf{Does not transfer:}
  \begin{itemize}
  \item Supermers and heavy-hitter handling, which rely on k-mer-counting
  semantics.
  \item Extraction-side staging and overlap, because our cost was a per-record
  seek.
  \item HySortK's \emph{static, size-based task dispatch}. We tested it on the
  alignment stage, which became the bottleneck after O1a. A task's cost there
  follows how much sequence it aligns, not its seed count: $R^2 = 0.19$ per
  task on primates4.chr6~\cite{pggbdata}. Greedy dispatch by seed count was
  no better than the existing length-based assignment, and at 8 ranks it was
  worse.
  \end{itemize}
\item \textbf{Transfers:} the \emph{task abstraction}, i.e.\ many independent
units of work, more than there are ranks. It frames the follow-on scheduling
work.
\end{itemize}

\section{Follow-on work (outside O1)}

\paragraph{Infrastructure (pull requests open or merged).}
\begin{itemize}
\item An \texttt{ALNtoPAF} time limit, with a non-zero exit on failure.
\item Per-stage instrumentation, plus freeing dead buffers early: peak memory
$-$25 to $-$38\%.
\item Exchanges above 2\,GB per rank, in bounded rounds. Rank outputs are
byte-identical with many rounds, and the seed exchange's peak dropped from
2.5\,GB to 40\,MB.
\item Scratch files can be placed off memory-backed \texttt{/tmp}.
\end{itemize}

\paragraph{Larger datasets.} Inputs from the PGGB data~\cite{pggb,pggbdata}
all match FastGA gold:
\begin{itemize}
\item scerevisiae8: 50,046 alignments;
\item primates4.chr6: 1,547,485 alignments, at 4--16 ranks;
\item primates14.chr6: 4,577,556 alignments, at 16 ranks and at 32 ranks
across 2 nodes.
\end{itemize}

\paragraph{Alignment scheduling (proposal O3).}
\begin{itemize}
\item On primates14.chr6 the alignment stage is 96--97\% of the run time.
The slowest rank does 3.35$\times$ the average at 32 ranks.
\item Self-alignment tasks are 49\% of the work.
\item Modelled from the measured task times, a dynamic task queue that runs
self-alignments of long contigs first would cut the slowest rank's alignment
time by 33--42\% at 16--32 ranks.
\item Beyond about 16 ranks, the single largest task (773\,s) is a floor for
any scheduler. Scaling further requires splitting self-alignment tasks while
keeping FastGA-identical output. That is an open question; prior distributed
aligners such as diBELLA~\cite{dibella} also leave work-aware alignment
assignment open.
\end{itemize}

\section{Conclusion}

O1 is complete. One targeted change, found by measurement rather than by
porting a technique, made extraction negligible: 1.5\,s of 65\,s at 1 rank and
0.2\,s of 10.6\,s at 8, with output identical to FastGA's.

The HySortK-inspired extensions planned for extraction were evaluated with
data and dropped. Of HySortK's architecture, the task abstraction carries
forward; size-based dispatch does not, because alignment cost depends on
output rather than input.

\bibliographystyle{ACM-Reference-Format}
\bibliography{o1-final-report}

\end{document}
